\lecture{7.6}{Use in Randomized Trials}
\label{lec:7-6}

\subsection{Why Covariate Adjustment Can Improve Precision}

In a randomized trial, treatment assignment is controlled by design. If treatment is
randomized with probability $p$, then
\[
A\ci W,\qquad \Prob(A=1\given W)=p.
\]
Because treatment is randomized, the unadjusted difference in means is consistent for
the ATE:
\[
\hat\psi_{\mathrm{unadj}}
=
\frac{\sum_iY_i\ind(A_i=1)}{\sum_i\ind(A_i=1)}
-
\frac{\sum_iY_i\ind(A_i=0)}{\sum_i\ind(A_i=0)}.
\]

Covariate adjustment is not required to remove confounding in a well-conducted
randomized trial. It can still improve precision. If baseline covariates strongly
predict the outcome, adjusting for them can remove outcome variation unrelated to
treatment, leading to smaller standard errors.

\subsection{Known Propensity Score Under Randomization}

In a trial with known randomization probability $p$,
\[
g_1(W)=p,\qquad g_0(W)=1-p.
\]
For equal randomization, $p=0.5$. In AIPTW, the known propensity score can be used
instead of estimating $g$:
\[
\psi_{n,\AIPTW}(1)
=
\frac{1}{n}\sum_{i=1}^n
\left[
\Qbar_n(1,W_i)
+
\frac{A_i}{p}\{Y_i-\Qbar_n(1,W_i)\}
\right].
\]
The corresponding expression for $a=0$ uses $(1-A_i)/(1-p)$.

\begin{assumptionbox}
In a randomized trial, exchangeability follows from the randomization design only for
the trial population and intervention actually randomized. Other issues, such as
nonadherence, loss to follow-up, interference, or changes in the target population,
require additional causal assumptions.
\end{assumptionbox}

\subsection{Efficiency Gains}

If $\Qbar_n$ predicts $Y$ well, the residuals $Y-\Qbar_n(A,W)$ are less variable than
the raw outcomes. AIPTW and TMLE use this predictive information while preserving
validness through the known randomization probability. The gain is largest when
baseline covariates are strongly prognostic.

For equal randomization and a well-estimated outcome regression, the precision gain
depends on how much variation in the outcome is explained by baseline covariates and
how treatment effects vary with covariates.

\begin{figure}[H]
\centering
\begin{tikzpicture}[
  node distance=1.35cm,
  box/.style={draw, rounded corners=2pt, minimum width=3.2cm, minimum height=.78cm},
  arr/.style={-{Latex[length=2.5mm]}, thick}
]
\node[box] (rand) {Randomization\\ known $g$};
\node[box, below left=1cm and .8cm of rand] (unadj) {Unadjusted\\ difference};
\node[box, below right=1cm and .8cm of rand] (adj) {AIPTW/TMLE\\ adjusted};
\node[box, below=1cm of adj] (prec) {Potentially smaller\\ variance};
\draw[arr] (rand) -- (unadj);
\draw[arr] (rand) -- (adj);
\draw[arr] (adj) -- (prec);
\end{tikzpicture}
\caption{In randomized trials, adjustment is used for precision rather than confounding control.}
\label{fig:ch7-trial-adjustment}
\end{figure}

\subsection{Outcome Regression in Randomized Trials}

The outcome regression in a trial can include baseline covariates and treatment:
\[
\Qbar(a,w)=\E(Y\given A=a,W=w).
\]
It may include treatment-covariate interactions if effect heterogeneity is plausible.
Because the propensity score is known, double robustness implies that consistency of
AIPTW does not require the outcome regression to be correct. However, a poor outcome
regression can fail to improve precision and may slightly harm finite-sample
performance.

\subsection{When Doubly Robust Methods Are Useful Even Without Confounding}

Doubly robust methods are useful in randomized trials because they provide:
\begin{itemize}
\item design-protected consistency through the known $g$;
\item valid influence-function inference under mild conditions;
\item possible efficiency gains from prognostic baseline covariates;
\item a unified analysis strategy that extends naturally to trials with unequal
randomization or covariate-adaptive designs.
\end{itemize}

\begin{keyidea}
In observational studies, adjustment is needed for identification. In randomized
trials, adjustment is often used for efficiency. The same AIPTW and TMLE formulas
apply, but the known randomization probability gives a particularly strong foundation.
\end{keyidea}
